\documentclass[11pt]{article}
\usepackage{amsmath,amssymb}
\usepackage[a4paper,margin=1in]{geometry}
\setlength{\parindent}{0pt}

\begin{document}

\title{A Concise Derivation of Classical Multidimensional Scaling}
\author{Zhenwei Tang}
\date{}
\maketitle

% ─────────────────────────────────────────────
\section*{Goal}
% ─────────────────────────────────────────────

Given $n$ objects with pairwise \emph{squared} dissimilarities
\[
  d_{ij}^{2} \in \mathbb{R}_{\ge 0}, \qquad i,j = 1,\dots,n,
\]
Classical MDS (cMDS) aims to find low-dimensional coordinate vectors
\[
  z_i \in \mathbb{R}^k, \qquad i = 1,\dots,n, \quad k \ll n,
\]
such that the Euclidean distances in the embedding space recover the
given dissimilarities as faithfully as possible:
\[
  \|z_i - z_j\|^2 \approx d_{ij}^{2}.
\]

The key insight of cMDS is to work with \emph{inner products} rather than
distances directly, because inner products admit a clean matrix
factorisation.  The derivation therefore proceeds in two phases:
(1) convert $d_{ij}^{2}$ into a matrix of inner products, and
(2) recover the coordinates via eigendecomposition.

% ─────────────────────────────────────────────
\section*{Step 1: From Distances to Inner Products}
% ─────────────────────────────────────────────

Suppose the (unknown) coordinates $z_1,\dots,z_n \in \mathbb{R}^k$ are
\emph{centred}:
\[
  \sum_{i=1}^{n} z_i = \mathbf{0}.          \tag{C}
\]
Define the \textbf{Gram matrix} $B \in \mathbb{R}^{n\times n}$ by
\[
  B_{ij} = z_i^{\top} z_j.
\]
Our goal is to express $B_{ij}$ purely in terms of $d_{ij}^{2}$, so that
we can compute $B$ from the given dissimilarities without knowing $z_i$.

Expand the squared distance:
\[
  d_{ij}^{2} = \|z_i - z_j\|^2
             = \|z_i\|^2 - 2\,z_i^{\top}z_j + \|z_j\|^2
             = B_{ii} - 2B_{ij} + B_{jj}.     \tag{1}
\]

Sum (1) over $j = 1,\dots,n$ and use the centring condition (C):
\[
  \sum_{j=1}^{n} d_{ij}^{2}
    = n\,B_{ii} - 2\sum_{j}B_{ij} + \sum_{j}B_{jj}
    = n\,B_{ii} + \operatorname{tr}(B),
\]
because $\sum_{j}B_{ij} = z_i^{\top}\!\sum_j z_j = 0$ by (C).  Hence
\[
  \frac{1}{n}\sum_{j=1}^{n} d_{ij}^{2}
    = B_{ii} + \frac{1}{n}\operatorname{tr}(B).  \tag{2}
\]

Similarly, summing (1) over $i$:
\[
  \frac{1}{n}\sum_{i=1}^{n} d_{ij}^{2}
    = B_{jj} + \frac{1}{n}\operatorname{tr}(B).  \tag{3}
\]

Summing (1) over both $i$ and $j$:
\[
  \frac{1}{n^2}\sum_{i,j} d_{ij}^{2}
    = \frac{2}{n}\operatorname{tr}(B).            \tag{4}
\]

Now solve for $B_{ij}$.  From (1):
\[
  B_{ij} = \tfrac{1}{2}\!\left(B_{ii}+B_{jj}-d_{ij}^{2}\right).
\]
Substitute (2), (3), and (4) to eliminate $B_{ii}$, $B_{jj}$, and
$\operatorname{tr}(B)$:

\[
  B_{ij}
  = -\tfrac{1}{2}
    \!\left(
      d_{ij}^{2}
      - \frac{1}{n}\sum_{j'}d_{ij'}^{2}
      - \frac{1}{n}\sum_{i'}d_{i'j}^{2}
      + \frac{1}{n^2}\sum_{i',j'}d_{i'j'}^{2}
    \right).   \tag{5}
\]

\textbf{Why this works:} equation (5) eliminates the unknown norms
$B_{ii}, B_{jj}$ and the trace by taking row, column, and grand means of
$d_{ij}^{2}$—a standard \emph{double-centring} operation.

% ─────────────────────────────────────────────
\section*{Step 2: Matrix Form of Double Centring}
% ─────────────────────────────────────────────

Let $D \in \mathbb{R}^{n\times n}$ denote the matrix of squared
dissimilarities, $D_{ij} = d_{ij}^{2}$, and let
\[
  H = I_n - \frac{1}{n}\mathbf{1}\mathbf{1}^{\top}
\]
be the \textbf{centring matrix}, where $\mathbf{1} = (1,\dots,1)^{\top}
\in \mathbb{R}^n$.  $H$ is symmetric and idempotent ($H^2=H$); it
projects out the mean of any vector it acts on.

\textbf{Applying the double-centring formula (5) in matrix form}
(see Appendix~A):
\[
  \boxed{B = -\tfrac{1}{2}\,H D H.}   \tag{6}
\]

This single expression computes $B$ directly from the observed
dissimilarity matrix $D$.

% ─────────────────────────────────────────────
\section*{Step 3: Eigendecomposition of $B$}
% ─────────────────────────────────────────────

Because $B = Z Z^{\top}$ where $Z = [z_1 \mid \cdots \mid z_n]^{\top}
\in \mathbb{R}^{n\times k}$, the matrix $B$ is \emph{symmetric positive
semidefinite} (PSD; see Appendix~B).

\textbf{Applying the eigendecomposition} (see Appendix~C) to $B$:
\[
  B = V \Lambda V^{\top},
\]
where
\[
  \Lambda = \operatorname{diag}(\lambda_1,\dots,\lambda_n),
  \quad \lambda_1 \ge \lambda_2 \ge \cdots \ge \lambda_n \ge 0,
\]
and $V = [v_1 \mid \cdots \mid v_n]$ is orthogonal ($V^{\top}V = I$).

To obtain a $k$-dimensional embedding, retain only the $k$ largest
eigenvalues.  Partition:
\[
  \Lambda_k = \operatorname{diag}(\lambda_1,\dots,\lambda_k) \in
  \mathbb{R}^{k\times k},
  \qquad
  V_k = [v_1 \mid \cdots \mid v_k] \in \mathbb{R}^{n\times k}.
\]

\textbf{Why truncate at $k$?}  Retaining the top $k$ eigenpairs gives the
best rank-$k$ approximation to $B$ in the Frobenius norm
(Eckart–Young theorem, Appendix~D), and hence the embedding that
minimises the reconstruction error in the inner-product sense.

% ─────────────────────────────────────────────
\section*{Step 4: Recover the Coordinates}
% ─────────────────────────────────────────────

We need $Z_k \in \mathbb{R}^{n\times k}$ such that
$Z_k Z_k^{\top} \approx B$.  From the truncated decomposition:
\[
  B \approx V_k \Lambda_k V_k^{\top}
          = \bigl(V_k \Lambda_k^{1/2}\bigr)
            \bigl(V_k \Lambda_k^{1/2}\bigr)^{\top}.
\]

Hence we set
\[
  \boxed{Z_k = V_k\,\Lambda_k^{1/2},}   \tag{7}
\]
where $\Lambda_k^{1/2} = \operatorname{diag}(\sqrt{\lambda_1},\dots,
\sqrt{\lambda_k})$.

The $i$-th row of $Z_k$ gives the $k$-dimensional embedding of object $i$:
\[
  z_i = \Lambda_k^{1/2} V_k^{\top} e_i \in \mathbb{R}^k,
\]
where $e_i$ is the $i$-th standard basis vector.

% ─────────────────────────────────────────────
\section*{Conclusion}
% ─────────────────────────────────────────────

Starting from pairwise squared dissimilarities $D_{ij} = d_{ij}^2$, cMDS
recovers low-dimensional coordinates in three steps:

\begin{enumerate}
  \item \textbf{Double-centre} $D$ to obtain the Gram matrix of inner
        products:
        \[B = -\tfrac{1}{2}\,HDH.\]
  \item \textbf{Eigendecompose} $B$; retain the $k$ largest eigenpairs
        $(\lambda_i, v_i)$.
  \item \textbf{Form coordinates}
        $Z_k = V_k \Lambda_k^{1/2}$,
        whose rows are the desired $k$-dimensional embeddings.
\end{enumerate}

\begin{center}
\fbox{\parbox{0.85\textwidth}{\centering
  cMDS embeds objects in $\mathbb{R}^k$ by matching pairwise Euclidean
  distances to the given dissimilarities via inner-product factorisation.
}}
\end{center}

% ─────────────────────────────────────────────
\appendix
\section{Double Centring in Matrix Form}
% ─────────────────────────────────────────────

Let $H = I - \frac{1}{n}\mathbf{1}\mathbf{1}^{\top}$.
For any matrix $A$, the $(i,j)$ entry of $HAH$ is:
\[
  (HAH)_{ij}
  = A_{ij}
  - \frac{1}{n}\sum_{j'}A_{ij'}
  - \frac{1}{n}\sum_{i'}A_{i'j}
  + \frac{1}{n^2}\sum_{i',j'}A_{i'j'}.
\]
Applying this with $A = D$ and multiplying by $-\tfrac{1}{2}$ recovers
exactly formula (5), confirming $B = -\tfrac{1}{2}HDH$.

% ─────────────────────────────────────────────
\section{Symmetric Positive Semidefinite (PSD) Matrices}
% ─────────────────────────────────────────────

A matrix $M \in \mathbb{R}^{n\times n}$ is \textbf{symmetric positive
semidefinite} if:
\begin{itemize}
  \item $M = M^{\top}$ (symmetry), and
  \item $x^{\top}Mx \ge 0$ for all $x \in \mathbb{R}^n$.
\end{itemize}
Equivalently, all eigenvalues of $M$ are non-negative.

$B = ZZ^{\top}$ is always PSD:
$x^{\top}Bx = x^{\top}ZZ^{\top}x = \|Z^{\top}x\|^2 \ge 0$.

If the given dissimilarities are \emph{not} Euclidean, some eigenvalues of
$B$ may be (slightly) negative due to noise; in practice these are set to
zero.

% ─────────────────────────────────────────────
\section{Spectral Theorem (Eigendecomposition of Symmetric Matrices)}
% ─────────────────────────────────────────────

\textbf{Theorem.}  Every real symmetric matrix $M \in \mathbb{R}^{n\times
n}$ can be written as
\[
  M = V \Lambda V^{\top},
\]
where $V \in \mathbb{R}^{n\times n}$ is orthogonal ($V^{\top}V = I$) and
$\Lambda = \operatorname{diag}(\lambda_1,\dots,\lambda_n)$ contains the
real eigenvalues of $M$.

For a PSD matrix, all $\lambda_i \ge 0$, so $\Lambda^{1/2}$ is
well-defined.

% ─────────────────────────────────────────────
\section{Eckart–Young Theorem (Best Low-Rank Approximation)}
% ─────────────────────────────────────────────

\textbf{Theorem.}  Let $M = V\Lambda V^{\top}$ be the eigendecomposition
of a symmetric PSD matrix, with eigenvalues ordered $\lambda_1 \ge \cdots
\ge \lambda_n \ge 0$.  Among all rank-$k$ matrices $\hat{M}$, the one
minimising $\|M - \hat{M}\|_F^2$ is
\[
  \hat{M} = V_k \Lambda_k V_k^{\top},
\]
where $V_k$ and $\Lambda_k$ retain only the top $k$ eigenpairs.

In cMDS this justifies choosing the $k$ largest eigenvalues: the resulting
embedding minimises the total squared error between $B$ and $Z_kZ_k^{\top}$.

\end{document}